%-------------------------------------------------------------------- Comente: primero en mover
\needspace{8\baselineskip}
\tema{Comente: primer movimiento}
\question \textbf{Comente.} En modelos de competencia imperfecta, siempre es preferible fijar primero el precio (o la cantidad) que el rival.
\needspace{5\baselineskip}
\begin{solution}
Falso. La decisión depende de la variable con la cual se compita. Si son cantidades, es decir, à la Cournot, efectivamente es preferible mover primero, pero si se compite en precios, à la Bertrand, se obtienen mayores beneficios no moviendo en primer lugar.

\end{solution}

%-------------------------------------------------------------------- Comente: asimétricas
\needspace{8\baselineskip}
\tema{Comente: firmas asimétricas}
\question \textbf{Comente.} En un mercado con dos firmas asimétricas, la firma ineficiente siempre prefiere competir bajo Cournot, mientras que la firma eficiente siempre preferirá Bertrand.
\needspace{5\baselineskip}
\begin{solution}
Bertrand representa un nivel de competencia mayor. Si la superioridad en costos de la firma eficiente respecto de la ineficiente es superior a cierto umbral, la firma eficiente preferirá un nivel de competencia mayor (Bertrand). Si la diferencia en costos está por debajo de cierto umbral, la firma eficiente preferirá competir bajo Cournot. Es verdadero que la firma ineficiente siempre preferirá competir bajo Cournot, pues compitiendo bajo Bertrand obtendrá beneficios iguales a cero. Por lo tanto, la afirmación es falsa en cuanto a la firma eficiente.

\end{solution}

%-------------------------------------------------------------------- Comente: comparar
\needspace{8\baselineskip}
\tema{Comente: Cournot vs Bertrand}
\question \textbf{Comente.} Compare las competencias à la Cournot y Bertrand. ¿Qué mercados son mejor explicados por cada modelo?
\needspace{5\baselineskip}
\begin{solution}
La característica más evidente es la variable que manejan las empresas para competir (cantidades en Cournot, precios en Bertrand). En Cournot la competencia es menos fuerte que en Bertrand.

Si la capacidad y la producción se pueden ajustar rápidamente, el modelo de Bertrand es más apropiado (software, seguros, banca). Si la capacidad no se puede ajustar en el corto plazo, el modelo de Cournot es más apropiado (cemento, autos, electricidad, \emph{commodities} en general).

\end{solution}
